% GERADO por scripts/analyse_feedback.py — NÃO EDITAR À MÃO.
% Correr o script outra vez depois de a recolha fechar; o texto acompanha os dados.
\subsection{Utilidade percebida dos alertas entregues}
\label{sec:av_feedback}

O canal passou a acompanhar cada alerta entregue com dois botões, \emph{útil} e \emph{não ajudou}, e a registar o voto de quem carrega. As regras de análise foram fixadas antes de existir um único voto e não foram alteradas depois: mínimo de vinte votos efetivos para reportar qualquer proporção, um voto por pessoa e por alerta com o último a substituir o anterior, salvaguarda que repete o cálculo sem o votante dominante quando este representa mais de quarenta por cento dos votos, sempre sujeita ao mesmo mínimo, e intervalos de confiança de Wilson, calculados sem correção por agrupamento.
Contaram apenas votos sobre alertas presentes no registo partilhado.

Foram registados 131 votos válidos entre 2026-09-02 e 2026-09-11, que correspondem a 91 votos efetivos de 3 pessoas distintas, sobre 62 alertas. 11 votos foram posteriormente alterados pela mesma pessoa, e apenas o último de cada par conta. Outros 29 repetem um voto anterior sem o alterar e contam uma só vez. Outros 5 votos, que não entram nos 131 acima, foram excluídos por não corresponderem a nenhum alerta do registo partilhado.

Dos 91 votos efetivos, 86 classificaram o alerta como útil, ou seja 95\%, com intervalo de confiança de Wilson a 95\% entre 88\% e 98\%. Este cálculo binomial não corrige a dependência entre votos da mesma pessoa ou sobre o mesmo alerta; a sua largura não representa toda a incerteza desta amostra.

Uma única pessoa representa 58\% dos votos efetivos, acima do limite de 40\% fixado no protocolo. Excluindo-a, restam 38 votos, dos quais 34 classificam o alerta como útil.
A proporção sem essa pessoa é de 89\%, com intervalo entre 76\% e 96\%. É esta a leitura a reter.

A independência entre quem constrói o sistema e quem o classifica é condição desta amostra: o autor não figura entre os votantes contabilizados. Os votos provêm de leitores do canal público, e o registo conserva de cada um apenas o identificador de conversa atribuído pela plataforma.

Estes votos constituem retorno observacional em contexto real e não substituem o estudo controlado que a Secção~\ref{sec:con_limitacoes} identifica como a lacuna de maior peso deste trabalho, e que permanece por realizar.

Quatro limitações acompanham este resultado e nenhuma delas se resolve com mais votos. A primeira é a autosseleção: vota quem quer, e quem considera um alerta indiferente tende a não carregar em nada, o que empurra a amostra para os extremos. A segunda é a ausência de contrafactual, uma vez que nenhum grupo recebeu a variação de preço sem a explicação que a acompanha; nada aqui atribui a utilidade à explicação em si. A terceira é a distância entre utilidade percebida e decisão melhor, que é a hipótese fundadora deste trabalho e permanece por testar. A quarta é o desconhecimento de quem vota: o canal é público e não se sabe se quem responde pertence ao público que a dissertação assume.
